\section{Límites proyectivos en álgebra homológica}
\label{section:0.13}
\label{section:0.13-fr}

\subsection{La condición de Mittag-Leffler}
\label{subsection:0.13.1}
\label{subsection:0.13.1-fr}

\begin{env}[13.1.1]
\label{0.13.1.1}
\label{0.13.1.1-fr}
Sea $\boldsymbol{C}$ una categoría abeliana en la que existan los productos
infinitos (axioma AB~3*) de T, 1.5).
Entonces existe el ínfimo de una familia de subobjetos de un objeto de
$\boldsymbol{C}$, y todo sistema proyectivo de objetos de $\boldsymbol{C}$ posee un
límite proyectivo; este último es un funtor exacto por la izquierda del
sistema proyectivo en cuestión (T, 1.8).
\oldpage[0\textsubscript{III}]{65}
Sea $(A_\alpha,f_{\alpha\beta})$ un sistema proyectivo de objetos de
$\boldsymbol{C}$ cuyo conjunto de índices $I$ es \emph{filtrante} a la derecha; sean $A=\varprojlim A_\alpha$
y, para cada $\alpha\in I$, sea
$f_\alpha:A\to A_\alpha$ el morfismo canónico.
Para todo $\alpha\in I$, los subobjetos
$f_{\alpha\beta}(A_\beta)$, para $\alpha\leq\beta$, forman una familia
filtrante decreciente de subobjetos de $A_\alpha$.
El subobjeto
$A'_\alpha=\inf_{\beta\geq\alpha}f_{\alpha\beta}(A_\beta)$
se llama el subobjeto de las \emph{imágenes universales} en $A_\alpha$; está claro que $f_\alpha(A)\subset A'_\alpha$ y $f_{\alpha\beta}(A'_\beta)\subset A'_\alpha$ para $\alpha\leq\beta$; por tanto
$(A'_\alpha,f_{\alpha\beta}|A'_\beta)$ es un sistema proyectivo y
$A=\varprojlim A'_\alpha$.
\end{env}

\begin{env}[13.1.2]
\label{0.13.1.2}
\label{0.13.1.2-fr}
Para un sistema proyectivo $(A_\alpha,f_{\alpha\beta})$ en $\boldsymbol{C}$,
se llama \emph{condición de Mittag-Leffler} a la siguiente:
\begin{itemize}
  \item[(ML)]
  \emph{Para todo índice $\alpha$, existe un $\beta\geq\alpha$ tal
  que, para todo $\gamma\geq\beta$,}
  $f_{\alpha\gamma}(A_\gamma)=f_{\alpha\beta}(A_\beta)$.
\end{itemize}

Está claro que la condición (ML) se cumple si todos los
$f_{\alpha\beta}$ son \emph{epimorfismos}, la condición (ML) se satisface. Recíprocamente, si se satisface (ML), y si para todo $\alpha\in I$ sea $A'_\alpha$ el subobjeto de las imágenes
universales en $A_\alpha$, la restricción de $f_{\alpha\beta}$ a $A'_\beta$ es un \emph{epimorfismo}
$A'_\beta\to A'_\alpha$ para $\alpha\leq\beta$: en efecto, si $\gamma\geq\beta$ es tal que $f_{\beta\delta}(A_\delta)=f_{\beta\gamma}(A_\gamma)$ para $\delta\geq\gamma$, se tiene
$A'_\beta=f_{\beta\gamma}(A_\gamma)$ y esto implica, por otra parte, $f_{\alpha\delta}(A_\delta)=f_{\alpha\gamma}(A_\gamma)$ para $\delta\geq\gamma$, por tanto $A'_\alpha=f_{\alpha\gamma}(A_\gamma)=f_{\alpha\beta}(A'_\beta)$.

Observemos también que la condición (ML) se satisface cuando los objetos $A_\alpha$ son
\emph{artinianos} en $\boldsymbol{C}$, es decir, cuando toda familia de
subobjetos de $A_\alpha$ posee un elemento minimal.
En efecto, un elemento minimal de la familia filtrante decreciente
$(f_{\alpha\beta}(A_\beta))$ de subobjetos de $A_\alpha$
es necesariamente el menor de esos subobjetos.
\end{env}

\begin{remark}[13.1.3]
\label{0.13.1.3}
\label{0.13.1.3-fr}
\emph{Observación} (13.1.3).
La condición (ML) puede formularse también cuando $\boldsymbol{C}$ es, por
ejemplo, la categoría de conjuntos.
Se puede definir de nuevo el subconjunto de las imágenes universales en
$A_\alpha$, y las observaciones hechas en \sref{0.13.1.1} y
\sref{0.13.1.2} siguen siendo válidas.
\end{remark}

\subsection{La condición de Mittag-Leffler para grupos abelianos}
\label{subsection:0.13.2}
\label{subsection:0.13.2-fr}

\begin{proposition}[13.2.1]
\label{0.13.2.1}
\label{0.13.2.1-fr}
Sea
\[
0\longrightarrow A_\alpha\xrightarrow{u_\alpha}B_\alpha
\xrightarrow{v_\alpha}C_\alpha\longrightarrow0
\]
una sucesión exacta de sistemas proyectivos de grupos abelianos, relativos
al mismo conjunto de índices filtrante $I$.
\begin{enumerate}
  \item[{\rm(i)}]
  Si $(B_\alpha)$ satisface {\rm(ML)}, entonces también $(C_\alpha)$.
  \item[{\rm(ii)}]
  Si $(A_\alpha)$ y $(C_\alpha)$ satisfacen {\rm(ML)}, entonces también
  $(B_\alpha)$.
\end{enumerate}
\end{proposition}

\begin{proof}
Sean $(f_{\alpha\beta})$, $(g_{\alpha\beta})$ y
$(h_{\alpha\beta})$ los sistemas de homomorfismos que definen
$(A_\alpha)$, $(B_\alpha)$ y $(C_\alpha)$ respectivamente.

(i) Supongamos que $g_{\alpha\beta}(B_\beta)=g_{\alpha\lambda}(B_\lambda)$ para $\lambda\geq\beta$; como $v_\beta$ y $v_\lambda$ son sobreyectivos,
$h_{\alpha\beta}(C_\beta)=v_\alpha(g_{\alpha\beta}(B_\beta))=v_\alpha(g_{\alpha\lambda}(B_\lambda))=h_{\alpha\lambda}(C_\lambda)$
para $\lambda\geq\beta$.

(ii) Sea $\alpha\in I$, y sea $\beta\geq\alpha$ un índice tal que para $\lambda\geq\beta$ se tenga $f_{\alpha\beta}(A_\beta)=f_{\alpha\lambda}(A_\lambda)$; sea, por otra parte, $\gamma\geq\beta$ un índice tal que, para $\lambda\geq\gamma$, se tenga $h_{\beta\gamma}(C_\gamma)=h_{\beta\lambda}(C_\lambda)$. Sea entonces $y_\alpha$ un elemento de $g_{\alpha\gamma}(B_\gamma)$; se tiene, pues,
$y_\alpha=g_{\alpha\gamma}(y_\gamma)$ con $y_\gamma\in B_\gamma$,
y pongamos $y_\beta=g_{\beta\gamma}(y_\gamma)$.
Entonces
$v_\beta(y_\beta)=h_{\beta\gamma}(v_\gamma(y_\gamma))$.
Para todo $\lambda\geq\gamma$, existe, por hipótesis, un
$y_\lambda\in B_\lambda$ tal que
$h_{\beta\gamma}(v_\gamma(y_\gamma))
=h_{\beta\lambda}(v_\lambda(y_\lambda))
=v_\beta(g_{\beta\lambda}(y_\lambda))$,
De aquí que
$v_\beta(y_\beta\setmin g_{\beta\lambda}(y_\lambda))=0$, y por tanto
existe un $x_\beta\in A_\beta$
\oldpage[0\textsubscript{III}]{66}
tal que
$y_\beta=g_{\beta\lambda}(y_\lambda)+u_\beta(x_\beta)$.
Se sigue que
$y_\alpha=g_{\alpha\lambda}(y_\lambda)+u_\alpha(f_{\alpha\beta}(x_\beta))$
Como $\lambda\geq\beta$, existe un $x_\lambda\in A_\lambda$ tal
que
$f_{\alpha\beta}(x_\beta)=f_{\alpha\lambda}(x_\lambda)$,
Por consiguiente,
$y_\alpha=g_{\alpha\lambda}(y_\lambda+u_\lambda(x_\lambda))\in g_{\alpha\lambda}(B_\lambda)$,
lo que demuestra la aserción.
\end{proof}

\begin{proposition}[13.2.2]
\label{0.13.2.2}
\label{0.13.2.2-fr}
Sea $I$ un conjunto ordenado filtrante que posee un subconjunto
cofinal numerable, y sea
\[
0\longrightarrow A_\alpha\xrightarrow{u_\alpha}B_\alpha
\xrightarrow{v_\alpha}C_\alpha\longrightarrow0
\]
una sucesión exacta de sistemas proyectivos de grupos abelianos indexados
por $I$.
Si $(A_\alpha)$ satisface {\rm(ML)}, entonces la sucesión
\[
0\longrightarrow\varprojlim A_\alpha
\longrightarrow\varprojlim B_\alpha
\longrightarrow\varprojlim C_\alpha\longrightarrow0
\]
es exacta.
\end{proposition}

\begin{proof}
Solo queda demostrar que el homomorfismo
$v=\varprojlim v_\alpha:\varprojlim B_\alpha\to\varprojlim C_\alpha$ es sobreyectivo. Sea $z=(z_\alpha)$ un elemento de $\varprojlim C_\alpha$, y pongamos
$E_\alpha=v_\alpha^{-1}(z_\alpha)$
Los $E_\alpha$ forman manifiestamente un sistema proyectivo de conjuntos no
vacíos mediante las restricciones de los homomorfismos
$g_{\alpha\beta}:B_\beta\to B_\alpha$.

Demostremos que este sistema proyectivo satisface (ML).
Identificando $A_\alpha$ con un subgrupo de $B_\alpha$ por medio de
$u_\alpha$, para todo $\alpha\in I$, un índice
$\beta\geq\alpha$ tal que $g_{\alpha\beta}(A_\beta)=g_{\alpha\lambda}(A_\lambda)$ para $\lambda\geq\beta$; mostremos que se tiene también $g_{\alpha\beta}(E_\beta)=g_{\alpha\lambda}(E_\lambda)$ para $\lambda\geq\beta$. En efecto, tomemos un $y_\lambda\in E_\lambda$ y pongamos $y_\beta=g_{\beta\lambda}(y_\lambda)$, $y_\alpha=g_{\alpha\lambda}(y_\lambda)$; sea $y'_\alpha\in g_{\alpha\beta}(E_\beta)$, de modo que
$y'_\alpha=g_{\alpha\beta}(y'_\beta)$ para algún
$y'_\beta\in E_\beta$.
Entonces $y'_\beta\setmin y_\beta=x_\beta\in A_\beta$, y por hipótesis existe
un $x_\lambda\in A_\lambda$ tal que
$g_{\alpha\beta}(x_\beta)=g_{\alpha\lambda}(x_\lambda)$
Así,
\[
y'_\alpha=g_{\alpha\beta}(y_\beta)+g_{\alpha\beta}(x_\beta)
=g_{\alpha\lambda}(y_\lambda)+g_{\alpha\lambda}(x_\lambda)
=g_{\alpha\lambda}(y_\lambda+x_\lambda)
\in g_{\alpha\lambda}(E_\lambda),
\] lo que demuestra nuestra aserción. Dicho esto, se sabe (Bourbaki, \emph{Top. gén.}, cap.~II, 3\textsuperscript{a} ed., §~3, teorema~1) que bajo las hipótesis hechas sobre $I$, un sistema proyectivo de conjuntos no vacíos que satisface (ML)
posee un límite proyectivo no vacío.
Existe, por tanto, un punto
$y=(y_\alpha)\in\varprojlim E_\alpha$.
Como $v_\alpha(y_\alpha)=z_\alpha$ para todo $\alpha$, se tiene
$z=v(y)$.
\end{proof}

\begin{proposition}[13.2.3]
\label{0.13.2.3}
\label{0.13.2.3-fr}
Bajo las hipótesis sobre $I$ de \sref{0.13.2.2}, sea
$(K_\alpha^\bullet)_{\alpha\in I}$ un sistema proyectivo
de complejos de grupos abelianos
$K_\alpha^\bullet=(K_\alpha^n)_{n\in\mathbb Z}$
cuyas diferenciales tienen grado $+1$.
Para todo $n$, existe un homomorfismo canónico funtorial
\begin{equation}
\label{0.13.2.3.1-fr}
\tag{13.2.3.1}
h_n:H^n(\varprojlim K_\alpha^\bullet)
\longrightarrow\varprojlim H^n(K_\alpha^\bullet).
\end{equation}
Si, para todo grado $n$, el sistema proyectivo de grupos abelianos
$(K_\alpha^n)_{\alpha\in I}$ satisface {\rm(ML)},
entonces todo $h_n$ son sobreyectivos. Si además, para un grado $n$, el sistema proyectivo $(H^{n-1}(K_\alpha^\bullet))_{\alpha\in I}$
satisface {\rm(ML)}, entonces $h_n$ es biyectivo.
\end{proposition}

\begin{proof}
Pongamos, para todo $n$, $K^n=\varprojlim K_\alpha^n$; la definición de los homomorfismos $h_n$ están definidos por los diagramas conmutativos
\[
\xymatrix{
\ldots\ar[r]&K^{n-1}\ar[r]\ar[d]&K^n\ar[r]\ar[d]
&K^{n+1}\ar[r]\ar[d]&\ldots\\
\ldots\ar[r]&K_\alpha^{n-1}\ar[r]&K_\alpha^n\ar[r]
&K_\alpha^{n+1}\ar[r]&\ldots
}
\]
\oldpage[0\textsubscript{III}]{67}
siendo las diferenciales en $K^\bullet$ los límites proyectivos
de las diferenciales correspondientes en los
$K_\alpha^\bullet$.
Consideremos las sucesiones exactas
\[
(*_n)\qquad
0\longrightarrow \mathrm B^n(K_\alpha^\bullet)
\longrightarrow \mathrm Z^n(K_\alpha^\bullet)
\longrightarrow H^n(K_\alpha^\bullet)\longrightarrow0
\]
y
\[
(**_n)\qquad
0\longrightarrow \mathrm Z^{n-1}(K_\alpha^\bullet)
\longrightarrow K_\alpha^{n-1}
\longrightarrow \mathrm B^n(K_\alpha^\bullet)\longrightarrow0
\]
La hipótesis y la Proposición~\sref{0.13.2.1}~(i) muestran que el
sistema proyectivo
$(\mathrm B^n(K_\alpha^\bullet))_{\alpha\in I}$
satisface (ML) para todo $n$.
Se sigue de \sref{0.13.2.2} que
\[
(***_n)\qquad
0\longrightarrow\varprojlim_\alpha\mathrm B^n(K_\alpha^\bullet)
\longrightarrow\varprojlim_\alpha\mathrm Z^n(K_\alpha^\bullet)
\longrightarrow\varprojlim_\alpha H^n(K_\alpha^\bullet)
\longrightarrow0
\]
es exacta.
Ahora bien,
$\varprojlim\mathrm B^n(K_\alpha^\bullet)$
se identifica con un subgrupo de $K^n$ que contiene
$\mathrm B^n(K^\bullet)$, mientras que
$\varprojlim\mathrm Z^n(K_\alpha^\bullet)$
se identifica con un subgrupo de $\mathrm Z^n(K^\bullet)$.
Por consiguiente, $h_n$ es sobreyectiva. Si ahora se supone además que el sistema proyectivo $(H^{n-1}(K_\alpha^\bullet))_{\alpha\in I}$
satisface (ML).
Las sucesiones exactas $(*_{n-1})$ y la
Proposición~\sref{0.13.2.1}~(ii) muestran que
$(\mathrm Z^{n-1}(K_\alpha^\bullet))_{\alpha\in I}$
satisface (ML).
Aplicando \sref{0.13.2.2} a las sucesiones exactas $(**_n)$ se obtiene una
sucesión exacta
\[
0\longrightarrow
\varprojlim_\alpha\mathrm Z^{n-1}(K_\alpha^\bullet)
\longrightarrow K^{n-1}\xrightarrow{u}
\varprojlim_\alpha\mathrm B^n(K_\alpha^\bullet)
\longrightarrow0
\]
Como
$\varprojlim\mathrm B^n(K_\alpha^\bullet)
\supset\mathrm B^n(K^\bullet)$, y que el compuesto de la inyección $\varprojlim\mathrm B^n(K_\alpha^\bullet)\to K^n$ y de $u$ es el operador de derivación $K^{n-1}\to K^n$, la sobreyectividad de
$u$ implica
$\varprojlim\mathrm B^n(K_\alpha^\bullet)
=\mathrm B^n(K^\bullet)$,
Así, $h_n$ es inyectivo.
\end{proof}

\begin{env}[13.2.4]
\label{0.13.2.4}
\label{0.13.2.4-fr}
\emph{Observaciones} (13.2.4).
\begin{enumerate}
  \item[{\rm(i)}]
  El argumento de \sref{0.13.2.2}
  (cf.\ Bourbaki, \emph{loc.\ cit.}) muestra que la conclusión de
  esa proposición sigue siendo válida si se supone únicamente que los
  $A_\alpha$ pueden dotarse de estructuras de espacios metrizables completos
  en los que las traslaciones son homeomorfismos, que las aplicaciones
  $f_{\alpha\beta}:A_\beta\to A_\alpha$ que definen el sistema proyectivo $(A_\alpha)$ son \emph{uniformemente continuas} para las distancias consideradas, y finalmente que el sistema $(A_\alpha)$ satisface la condición
  \begin{itemize}
    \item[(ML$'$)]
    \emph{Para todo índice $\alpha$, existe un $\beta\geq\alpha$ tal
    que, para todo $\gamma\geq\beta$,
    $f_{\alpha\gamma}(A_\gamma)$ es denso en
    $f_{\alpha\beta}(A_\beta)$.}
  \end{itemize}

  Esto da un suplemento análogo a \sref{0.13.2.3}.
  Supongamos que $K_\alpha^n=0$ para $n<0$ y todo $\alpha$; supongamos además que $(K_\alpha^n)_{\alpha\in I}$ satisface (ML) para
  $n\geq0$; y que
  $A_\alpha=H^0(K_\alpha^\bullet)$ puede dotarse de
  estructuras de espacio métrico que tengan las propiedades recién enunciadas.
  Entonces las conclusiones de \sref{0.13.2.3} no cambian para
  $n\geq2$, y $h_1$ es además \emph{biyectivo}.
  En efecto, el argumento de \sref{0.13.2.2} muestra también que
  $(\mathrm B^1(K_\alpha^\bullet))_{\alpha\in I}$
  satisface (ML), que la sucesión $(***_1)$ es exacta, y finalmente, en virtud de lo que precede, que $\varprojlim\mathrm B^1(K_\alpha^\bullet)=\mathrm B^1(K^\bullet)$.
  Esto demuestra, entre otras cosas, las aserciones de (T, 3.10.2).

  \item[{\rm(ii)}]
  Se pueden introducir los funtores derivados por la derecha
  $\varprojlim^{(n)}$ del funtor $\varprojlim$, y obtener enunciados más
  completos que los precedentes~\cite{III-28}.
\end{enumerate}
\end{env}

\subsection{Aplicación: cohomología de un límite proyectivo de haces}
\label{subsection:0.13.3}
\label{subsection:0.13.3-fr}

\begin{proposition}[13.3.1]
\label{0.13.3.1}
\label{0.13.3.1-fr}
\oldpage[0\textsubscript{III}]{68}
Sea $X$ un espacio topológico, sea
$(\mathcal F_k)_{k\in\mathbb N}$ un sistema proyectivo de haces de
grupos abelianos sobre $X$, y sea $\mathcal F=\varprojlim\mathcal F_k$.
Se suponen verificadas las condiciones siguientes:
\begin{enumerate}
  \item[{\rm(i)}]
  Existe una base $\mathfrak B$ de la topología de $X$ tal que,
  para todo $U\in\mathfrak B$ y todo $i\geq0$, el sistema
  proyectivo
  $(H^i(U,\mathcal F_k))_{k\in\mathbb N}$ satisface \textup{(ML)}.

  \item[{\rm(ii)}]
  Para todo $x\in X$ y todo $i>0$, se tiene
  $\varinjlim_U\bigl(\varprojlim_k H^i(U,\mathcal F_k)\bigr)=0$
  donde $U$ recorre las vecindades de $x$ pertenecientes a
  $\mathfrak B$.

  \item[{\rm(iii)}]
  Los homomorfismos
  $u_{hk}:\mathcal F_k\to\mathcal F_h$ $(h\leq k)$ que definen el
  sistema proyectivo $(\mathcal F_k)$ son sobreyectivos.
\end{enumerate}
Bajo estas condiciones, para todo $i>0$, el homomorfismo canónico
\[
h_i:H^i(X,\mathcal F)\longrightarrow
\varprojlim_k H^i(X,\mathcal F_k)
\]
es sobreyectivo. Si, además, para un valor dado de $i$ el sistema
proyectivo
$(H^{i-1}(X,\mathcal F_k))_{k\in\mathbb N}$ satisface \textup{(ML)}, entonces $h_i$ es
biyectivo.
\end{proposition}

\begin{proof}
\emph{a)}
Supongamos primero que los $\mathcal F_k$, así como los \emph{núcleos}
$\mathcal N_{hk}$ de los $u_{hk}$, son \emph{flácidos}; mostraremos entonces que
la condición~{\rm(iii)} del enunciado implica
$H^i(X,\mathcal F)=0$ para $i>0$.
Basta demostrar que, para \emph{todo} subconjunto abierto $U$ de $X$ y todo
recubrimiento $\mathfrak U$ de $U$ por subconjuntos abiertos de $U$, se tiene $H^i(\mathfrak U,\mathcal F)=0$ para $i>0$. De ello resultará en efecto, en primer lugar, que para la cohomología de Čech se tiene
$\check H^i(U,\mathcal F)=0$ para todo $i>0$, y por tanto, por
(G, II, 5.9.2) aplicado al conjunto de todos los subconjuntos abiertos de
$X$, que $H^i(X,\mathcal F)=0$ para todo $i>0$.

Como los $\mathcal F_k$ son flácidos, se tiene
$H^i(\mathfrak U,\mathcal F_k)=0$ para todo $i>0$ (G, II, 5.2.3); consideremos para cada $k$ el complejo $C^\bullet(\mathfrak U,\mathcal F_k)$ de cocadenas
alternadas sobre el nervio del recubrimiento $\mathfrak U$ (G, II, 5.1), que forman evidentemente un sistema proyectivo de complejos de grupos abelianos. Mostremos que todas las aplicaciones $C^\bullet(\mathfrak U,\mathcal F_k)\to
C^\bullet(\mathfrak U,\mathcal F_h)$ $(h\leq k)$ son \emph{sobreyectivas}. Evidentemente, basta, por definición, demostrar que para todo abierto $V$
de $X$, la aplicación
$\Gamma(V,\mathcal F_k)\to\Gamma(V,\mathcal F_h)$ es
es sobreyectiva. Pero la sucesión
$0\to\mathcal N_{hk}\to\mathcal F_k\to\mathcal F_h\to0$
es exacta por hipótesis y da, por tanto, la sucesión exacta de
cohomología
\[
\Gamma(V,\mathcal F_k)\longrightarrow\Gamma(V,\mathcal F_h)
\longrightarrow H^1(V,\mathcal N_{hk})=0
\]
pues $\mathcal N_{hk}$ es flácido.

El sistema proyectivo
$(C^\bullet(\mathfrak U,\mathcal F_k))_{k\in\mathbb N}$
satisface, por tanto, (ML), y lo mismo vale para
$(H^i(\mathfrak U,\mathcal F_k))_{k\in\mathbb N}$ para todo $i\geq0$:
para $i>0$, y como $H^0(\mathfrak U,\mathcal F_k)=\Gamma(U,\mathcal F_k)$
(G, II, 5.2.2), la condición (ML) se cumple también para $i=0$ por lo que precede. Se puede, pues, aplicar (\hyperref[0.13.2.3-fr]{13.2.3}), que muestra que $H^i(\mathfrak U,\mathcal F)=\varprojlim_k H^i(\mathfrak U,\mathcal F_k)=0$ para todo $i>0$.

\emph{b)} Pasemos al caso general, y consideremos para cada $k\in\mathbb N$, la \emph{resolución canónica} $\mathcal C^\bullet(X,\mathcal F_k)
=(\mathcal C^i(X,\mathcal F_k))_{i\geq0}$
de $\mathcal F_k$ por haces flácidos (G, II, 4.3).
Para todo $i\geq0$, está claro que
$(\mathcal C^i(X,\mathcal F_k))_{k\in\mathbb N}$ es un sistema proyectivo de
haces flácidos; demostremos que satisface las condiciones de la
parte~\emph{a)}.
\oldpage[0\textsubscript{III}]{69}
En efecto, si $\mathcal N_{hk}$ denota el núcleo de $u_{hk}$ para $h\leq k$,
entonces la sucesión
$0\to\mathcal N_{hk}\to\mathcal F_k\to\mathcal F_h\to0$
es exacta por~{\rm(iii)}, y nuestra aserción se sigue de la exactitud
del funtor
$\mathcal A\leadsto\mathcal C^i(X,\mathcal A)$ es exacto (G, II, 4.3). Sea $\mathcal G^i=\varprojlim_k\mathcal C^i(X,\mathcal F_k)$
Por la parte~\emph{a)}, se tiene
$H^j(X,\mathcal G^i)=0$ para $j>0$ e $i\geq0$.
Demostraremos que
$\mathcal G^\bullet=(\mathcal G^i)_{i\geq0}$ es una \emph{resolución} del haz
$\mathcal F$. Como $H^j(X,\mathcal G^\bullet)=0$ para $j>0$, la cohomología
$H^\bullet(X,\mathcal F)$ será entonces igual a
$H^\bullet(\Gamma(X,\mathcal G^\bullet))$ (G, II, 4.7.1). Está claro que, por paso al límite proyectivo, se deducen de las sucesiones exactas \[
0\longrightarrow\mathcal F_k
\longrightarrow\mathcal C^0(X,\mathcal F_k)
\longrightarrow\mathcal C^1(X,\mathcal F_k)\longrightarrow\ldots
\]
se obtiene manifiestamente un complejo de haces de grupos abelianos
\[
0\longrightarrow\mathcal F\longrightarrow\mathcal G^0
\longrightarrow\mathcal G^1\longrightarrow\ldots
\]
Para demostrar nuestra aserción, queda establecer que
$\mathcal H^i(\mathcal G^\bullet)=0$ para $i>0$.
Este haz está generado por el prehaz
$U\leadsto H^i(\Gamma(U,\mathcal G^\bullet))$
(G, II, 4.1).
Ahora bien, el complejo $\Gamma(U,\mathcal G^\bullet)$ es el límite proyectivo
del sistema proyectivo de complejos de grupos abelianos
$(\Gamma(U,\mathcal C^\bullet(X,\mathcal F_k)))_{k\in\mathbb N}$
(\textbf{0}\textsubscript{I}, \EGAUnavailableHyperref{0.3.2.6}{3.2.6}).
Se ha visto en la parte~\emph{a)} que, para todo $i\geq0$, las aplicaciones $\Gamma(U,\mathcal C^i(X,\mathcal F_k))\to
\Gamma(U,\mathcal C^i(X,\mathcal F_h))$ $(h\leq k)$ son \emph{sobreyectivas}; por otra parte, se tiene
$H^i(U,\mathcal F_k)=
H^i(\Gamma(U,\mathcal C^\bullet(X,\mathcal F_k)))$,
pues la resolución canónica $\mathcal C^\bullet(U,\mathcal F_k|U)$ es
inducida sobre $U$ por $\mathcal C^\bullet(X,\mathcal F_k)$.
Por la hipótesis~{\rm(i)}, para todo $U\in\mathfrak B$, se puede aplicar (\hyperref[0.13.2.3-fr]{13.2.3}) al sistema proyectivo de complejos $(\Gamma(U,\mathcal C^\bullet(X,\mathcal F_k)))_{k\in\mathbb N}$, y se tiene, pues, $H^i(\Gamma(U,\mathcal G^\bullet))=\varprojlim_k H^i(U,\mathcal F_k)$ para todo $i\geq0$. La hipótesis (ii) demuestra entonces, por definición, que los haces $\mathcal H^i(\mathcal G^\bullet)$ se anulan para $i>0$.

Así, para todo $i\geq0$, $H^i(X,\mathcal F)=H^i(\Gamma(X,\mathcal G^\bullet))$ y \[
\Gamma(X,\mathcal G^\bullet)
=\varprojlim_k\Gamma(X,\mathcal C^\bullet(X,\mathcal F_k)).
\] Se acaba de observar que las aplicaciones $\Gamma(X,\mathcal C^i(X,\mathcal F_k))\to
\Gamma(X,\mathcal C^i(X,\mathcal F_h))$ $(h\leq k)$ son todas sobreyectivas; la conclusión resulta, pues, de nuevo de (\hyperref[0.13.2.3-fr]{13.2.3}).
\end{proof}

\begin{env}[13.3.2]
\label{0.13.3.2}
\label{0.13.3.2-fr}
\emph{Observaciones} (13.3.2).
\begin{enumerate}
  \item[{\rm(i)}]
  El enunciado de \sref{0.13.3.1} solo tiene interés para $i>0$,
  pues, para $i=0$, $h_i$ es siempre un isomorfismo sin hipótesis
  alguna (\textbf{0}\textsubscript{I}, \EGAUnavailableHyperref{0.3.2.6}{3.2.6}).

  \item[{\rm(ii)}]
  Las condiciones~{\rm(i)} y~{\rm(ii)} de \sref{0.13.3.1} se satisfacen,
  en particular, si
  $H^i(U,\mathcal F_k)=0$ para todo $k$, todo $i>0$ y todo
  $U\in\mathfrak B$, y si, para $U\in\mathfrak B$, las aplicaciones
  $\Gamma(U,\mathcal F_k)\to\Gamma(U,\mathcal F_h)$
  son \emph{sobreyectivas}. Este es el caso más frecuente en el que se aplica
  \sref{0.13.3.1}.
\end{enumerate}
\end{env}

\subsection{Condición de Mittag-Leffler y objetos graduados asociados a sistemas proyectivos}
\label{subsection:0.13.4}
\label{subsection:0.13.4-fr}
\label{0.13.4}
\label{0.13.4-fr}

\begin{env}[13.4.1]
\label{0.13.4.1}
\label{0.13.4.1-fr}
Sea
$\mathbf A=(A_k,u_{kh})_{k\in\mathbb Z}$ un sistema proyectivo en una
categoría abeliana $\boldsymbol{C}$.
Decimos que está \emph{acotado inferiormente} si existe un $k_0$ tal que
$A_k=0$ para $k<k_0$. Definiremos sobre cada $A_k$ una \emph{filtración} $(\mathrm F^p(A_k))_{p\in\mathbb Z}$ por
\begin{equation}
\label{0.13.4.1.1-fr}
\tag{13.4.1.1}
\begin{cases}
\mathrm F^p(A_k)=\operatorname{Ker}(A_k\longrightarrow A_{p-1})&\text{para }p\leq k+1\\
\mathrm F^p(A_k)=0&\text{para }p\geq k+1
\end{cases}
\end{equation}

\oldpage[0\textsubscript{III}]{70}
Por hipótesis,
$\mathrm F^{k_0}(A_k)=A_k$ y
$\mathrm F^{k+1}(A_k)=0$; en otras palabras, la filtración considerada
es finita (\sref{0.11.1.3}).
Los objetos graduados asociados a esta filtración son, por tanto,
\[
\operatorname{gr}^p(A_k)=
\operatorname{Ker}(A_k\longrightarrow A_{p-1})
/\operatorname{Ker}(A_k\longrightarrow A_p)
\]
Por consiguiente, $\operatorname{gr}^p(A_k)$ es isomorfo a la imagen por
$A_k\to A_p$ de
$\operatorname{Ker}(A_k\to A_{p-1})$.
Por la transitividad de los morfismos que definen un sistema proyectivo,
se tiene, por tanto,
\begin{equation}
\label{0.13.4.1.2-fr}
\tag{13.4.1.2}
\operatorname{gr}^p(A_k)=
\operatorname{Ker}(A_p\longrightarrow A_{p-1})
\cap\operatorname{Im}(A_k\longrightarrow A_p)
\end{equation}
Pero \sref{0.13.4.1.1} da
$\operatorname{Ker}(A_p\to A_{p-1})=\operatorname{gr}^p(A_p)$, y por tanto
\begin{equation}
\label{0.13.4.1.3-fr}
\tag{13.4.1.3}
\operatorname{gr}^p(A_k)=
\operatorname{gr}^p(A_p)\cap\operatorname{Im}(A_k\longrightarrow A_p).
\end{equation}
Las definiciones precedentes muestran además que se tiene para $k\leq h$,
\[
u_{kh}(\mathrm F^p(A_h))\subset\mathrm F^p(A_k)
\]
y por consiguiente que los $\operatorname{gr}^p(u_{kh})$ definen un sistema proyectivo
$(\operatorname{gr}^p(A_k))_{k\in\mathbb Z}$ para todo $p\in\mathbb Z$.
\end{env}

\begin{env}[13.4.2]
\label{0.13.4.2}
\label{0.13.4.2-fr}
Decimos que el sistema proyectivo $\mathbf A$ es
\emph{esencialmente constante} si los morfismos
$A_{k+1}\to A_k$ son isomorfismos para todo $k$
suficientemente grande.
Decimos que $\mathbf A$ es \emph{estricto} si los morfismos
$A_i\to A_j$ $(j\leq i)$ son epimorfismos.

Cuando $\mathbf A$ es estricto, se sigue de \sref{0.13.4.1.3} que, para
$p\leq k\leq h$, el morfismo canónico
$\operatorname{gr}^p(A_h)\to\operatorname{gr}^p(A_k)$
es un isomorfismo; en otras palabras, el sistema proyectivo
$(\operatorname{gr}^p(A_k))_{k\in\mathbb Z}$ es esencialmente constante.
La sucesión de objetos $\operatorname{gr}^p(A_p)$ (identificados con $\varprojlim_k\operatorname{gr}^p(A_k)$ para todo $p$) se denota entonces
$\operatorname{gr}^\bullet(\mathbf A)$ y se llama el
\emph{objeto graduado asociado al sistema proyectivo estricto}
$\mathbf A=(A_k)$. Si se supone ahora que el sistema proyectivo $\mathbf A$, acotado
inferiormente, satisface (ML).
Por \sref{0.13.1.2}, el sistema proyectivo
$\mathbf A'=(A'_k)$ de objetos de ``imagen universal'' es estricto; también está
acotado inferiormente.
El objeto graduado $\operatorname{gr}^\bullet(\mathbf A')$ asociado a
$\mathbf A'$ se llama asimismo el objeto graduado
\emph{asociado a} $\mathbf A$ y se denota
$\operatorname{gr}^\bullet(\mathbf A)$.
\end{env}

\begin{proposition}[13.4.3]
\label{0.13.4.3}
\label{0.13.4.3-fr}
Sea $\mathbf A=(A_k)_{k\in\mathbb Z}$ un sistema proyectivo acotado
inferiormente en una categoría abeliana $\boldsymbol{C}$.
Las condiciones siguientes son equivalentes:
\begin{enumerate}
  \item[{\rm a)}] $\mathbf A$ satisface \textup{(ML)}.
  \item[{\rm b)}] Para todo $p\in\mathbb Z$, el sistema proyectivo
  $(\operatorname{gr}^p(A_k))_{k\in\mathbb Z}$ es esencialmente constante.
\end{enumerate}
Además, cuando se cumplen estas condiciones, existe, para todo
$p\in\mathbb Z$, un isomorfismo canónico
\begin{equation}
\label{0.13.4.3.1-fr}
\tag{13.4.3.1}
\operatorname{gr}^p(\mathbf A)\xrightarrow{\sim}
\varprojlim_k\operatorname{gr}^p(A_k).
\end{equation}
\end{proposition}

\begin{proof}
Se sigue inmediatamente de \sref{0.13.4.1.2} que \emph{a)} implica
\emph{b)}. La misma fórmula, aplicada al sistema proyectivo
$\mathbf A'$ con la notación de \sref{0.13.4.2}, da
\sref{0.13.4.3.1} por definición.

Para $k\leq h$, pongamos
$A_{kh}=\operatorname{Im}(A_h\to A_k)$
Si $k\leq h\leq j$, entonces
$A_{kj}\subset A_{kh}\subset A_k$.
Dotemos a $A_{kh}$ de la filtración inducida por
$(\mathrm F^p(A_k))$.
La transitividad de los morfismos que definen $\mathbf A$ muestra
inmediatamente que esta es también la filtración cociente de
$(\mathrm F^p(A_h))$; por consiguiente,
\begin{equation}
\label{0.13.4.3.2-fr}
\tag{13.4.3.2}
\operatorname{gr}^p(A_{kh})=
\operatorname{Im}\bigl(\operatorname{gr}^p(A_h)\longrightarrow
\operatorname{gr}^p(A_k)\bigr).
\end{equation}

\oldpage[0\textsubscript{III}]{71}
Supongamos ahora que se cumple \emph{b)}.
Para todo $p\in\mathbb Z$ y todo $k\geq p$, existe un entero
$\mathrm L(p,k)$ tal que el miembro derecho de
\sref{0.13.4.3.2} es constante para $h\geq\mathrm L(p,k)$.
Como $\operatorname{gr}^p(A_k)=0$ para $p<k_0$, para $k$ fijo solo
un número finito de los enteros $\mathrm L(p,k)$ puede ser no nulo cuando
$p$ recorre los enteros $\leq k$.
Sea $\mathrm L(k)=m$ el mayor de estos enteros.
Para todo $h\geq m$, se tiene
$A_{kh}\subset A_{km}$, y, por la definición de $m$, la inyección
canónica
$A_{kh}\to A_{km}$ induce un isomorfismo
$\operatorname{gr}^\bullet(A_{kh})\xrightarrow{\sim}
\operatorname{gr}^\bullet(A_{km})$
Como las filtraciones son finitas, esta inyección es ella misma biyectiva
(Bourbaki, \emph{Alg.\ comm.}, cap.~III, \textsection2,
n\textsuperscript{o}~8, teorema~1).
Esto demuestra que $\mathbf A$ satisface (ML).
\end{proof}

\begin{env}[13.4.4]
\label{0.13.4.4-fr}
Supongamos que en $\boldsymbol{C}$ el límite proyectivo $A=\varprojlim A_k$ existe. En las definiciones de (\hyperref[0.13.4.1-fr]{13.4.1}), se puede entonces reemplazar $A_k$
por $A$, y la filtración así definida sobre $A$ sigue satisfaciendo
\begin{equation}
\label{0.13.4.4.1-fr}
\tag{13.4.4.1}
\operatorname{gr}^p(A)=\operatorname{gr}^p(A_p)
\cap\operatorname{Im}(A\longrightarrow A_p).
\end{equation}
\end{env}

\begin{corollary}[13.4.5]
\label{0.13.4.5}
\label{0.13.4.5-fr}
Supongamos que $\boldsymbol{C}$ es la categoría de grupos abelianos.
Si el sistema proyectivo $\mathbf A$ satisface \textup{(ML)} y
$A=\varprojlim_k A_k$, entonces para todo $p\in\mathbb Z$ existe un
isomorfismo canónico
\begin{equation}
\label{0.13.4.5.1-fr}
\tag{13.4.5.1}
\operatorname{gr}^p(A)\xrightarrow{\sim}
\varprojlim_k\operatorname{gr}^p(A_k).
\end{equation}
\end{corollary}

\begin{proof}
En efecto,
$\operatorname{Im}(A_k\to A_p)=\operatorname{Im}(A\to A_p)$
para todo $k$ suficientemente grande
(Bourbaki, \emph{Top.\ g\'{e}n.}, cap.~II,
3\textsuperscript{e}~\'{e}d., \textsection3,
n\textsuperscript{o}~5, teorema~1), y la conclusión se sigue de
\sref{0.13.4.1.3} y \sref{0.13.4.4.1}.
\end{proof}

\subsection{Límites proyectivos de sucesiones espectrales de complejos filtrados}
\label{subsection:0.13.5}
\label{subsection:0.13.5-fr}

\begin{env}[13.5.1]
\label{0.13.5.1}
\label{0.13.5.1-fr}
Sea $\boldsymbol{C}$ una categoría abeliana y sea $X^\bullet$ un complejo
de objetos de $\boldsymbol{C}$ dotado de una filtración
$(\mathrm F^p(X^\bullet))_{p\in\mathbb Z}$ tal que
$\mathrm F^{p_0}(X^\bullet)=X^\bullet$ para algún índice $p_0$. Consideremos para cada $k\in\mathbb Z$ el complejo $X_k^\bullet=X^\bullet/\mathrm F^{k+1}(X^\bullet)$; está provisto canónicamente de la filtración formada por los $\mathrm F^p(X_k^\bullet)=
\mathrm F^p(X^\bullet)/\mathrm F^{k+1}(X^\bullet)$ para $p\leq k$ y por los $\mathrm F^p(X_k^\bullet)=0$ para $p\geq k+1$.
Además, existen morfismos canónicos
$X_{k+1}^\bullet\to X_k^\bullet$, que hacen de
$\mathbf X^\bullet=(X_k^\bullet)_{k\in\mathbb Z}$ un \emph{sistema proyectivo
de complejos filtrados de objetos de $\boldsymbol{C}$}.
Se observará que este sistema proyectivo es \emph{estricto} y que
$X_k^\bullet=0$ para $k<p_0$.
\end{env}

\begin{env}[13.5.2]
\label{0.13.5.2}
\label{0.13.5.2-fr}
Más generalmente, consideremos un sistema proyectivo estricto acotado
inferiormente
$\mathbf X^\bullet=(X_k^\bullet)_{k\in\mathbb Z}$ de complejos de objetos de
$\boldsymbol{C}$, y dotemos cada $X_k^\bullet$ de la filtración definida en
\sref{0.13.4.1} (aplicada en la categoría abeliana de complejos acotados
inferiormente de objetos de $\boldsymbol{C}$).
Los morfismos
$X_k^\bullet\to X_p^\bullet$ ($(p\leq k)$) se convierten entonces
en morfismos de complejos \emph{filtrados} con filtraciones finitas.
La funtorialidad de las sucesiones espectrales de complejos filtrados
(\sref{0.11.2.3}) muestra que los morfismos de transición de
$\mathbf X^\bullet$ inducen morfismos que hacen de
$\mathrm E(\mathbf X^\bullet)=(\mathrm E(X_k^\bullet))$
un \emph{sistema proyectivo de sucesiones espectrales}.
\end{env}

\begin{lemma}[13.5.3]
\label{0.13.5.3}
\label{0.13.5.3-fr}
Supongamos que el sistema proyectivo
$\mathbf X^\bullet=(X_k^\bullet)_{k\in\mathbb Z}$ de complejos filtrados se
obtiene como en \sref{0.13.5.2}. Entonces:
\begin{enumerate}
  \item[{\rm a)}] Para $r\geq p-p_0$, se tiene
  $\mathrm B_r^{pq}(X_k^\bullet)=\mathrm B_\infty^{pq}(X_k^\bullet)$
  para todo $k\in\mathbb Z$.
  \item[{\rm b)}] Para $k+1\leq p+r$, se tiene
  $\mathrm Z_r^{pq}(X_k^\bullet)=\mathrm Z_\infty^{pq}(X_k^\bullet)$.
  \item[{\rm c)}] Para $k+1\geq p+r$, los morfismos $\mathrm Z_r^{pq}(X_h^\bullet)\to\mathrm Z_r^{pq}(X_k^\bullet)$ y $\mathrm B_r^{pq}(X_h^\bullet)\to\mathrm B_r^{pq}(X_k^\bullet)$ son isomorfismos para todo $h\geq k$.
\end{enumerate}
\end{lemma}

\oldpage[0\textsubscript{III}]{72}
Estas tres propiedades se siguen inmediatamente de las definiciones de
\sref{0.11.2.2}, teniendo en cuenta que
$\mathrm F^{p-r+1}(X_k^\bullet)=X_k^\bullet$ cuando $p-r<p_0$.

\begin{env}[13.5.4]
\label{0.13.5.4-fr}
Supongamos satisfechas las hipótesis de (\hyperref[0.13.5.3-fr]{13.5.3}). Entonces, para $p,q,r$ fijos ($r$ finito), los sistemas proyectivos $(\mathrm Z_r^{pq}(X_k^\bullet))_{k\in\mathbb Z}$, $(\mathrm B_r^{pq}(X_k^\bullet))_{k\in\mathbb Z}$, $(\mathrm E_r^{pq}(X_k^\bullet))_{k\in\mathbb Z}$ son \emph{esencialmente constantes}; se denotará por $\mathrm Z_r^{pq}(\mathbf X^\bullet)$,
$\mathrm B_r^{pq}(\mathbf X^\bullet)$, y
$\mathrm E_r^{pq}(\mathbf X^\bullet)=
\mathrm Z_r^{pq}(\mathbf X^\bullet)/
\mathrm B_r^{pq}(\mathbf X^\bullet)$
Los objetos $\mathrm Z_r^{pq}(\mathbf X^\bullet)$ y
$\mathrm B_r^{pq}(\mathbf X^\bullet)$ se identifican canónicamente con
subobjetos de $\mathrm E_1^{pq}(\mathbf X^\bullet)$.
La definición de los $d_r^{pq}$ (M, XV, 1) muestra que estos
morfismos, tomados respecto de los $X_k^\bullet$, son también esencialmente
constantes, y definen, por tanto, morfismos
\begin{equation}
\label{0.13.5.4.1-fr}
\tag{13.5.4.1}
d_r^{pq}:\mathrm E_r^{pq}(\mathbf X^\bullet)
\longrightarrow\mathrm E_r^{p+r,q-r+1}(\mathbf X^\bullet)
\end{equation}
tal que
$d_r^{p+r,q-r+1}\circ d_r^{pq}=0$.
Además, existen isomorfismos canónicos de
$\operatorname{Ker}(d_r^{pq})$ sobre
$\mathrm Z_{r+1}^{pq}(\mathbf X^\bullet)/
\mathrm B_r^{pq}(\mathbf X^\bullet)$
y de $\operatorname{Im}(d_r^{pq})$ sobre
$\mathrm B_{r+1}^{p+r,q-r+1}(\mathbf X^\bullet)/
\mathrm B_r^{p+r,q-r+1}(\mathbf X^\bullet)$.
\end{env}

\begin{lemma}[13.5.5]
\label{0.13.5.5}
\label{0.13.5.5-fr}
Bajo las hipótesis de \sref{0.13.5.3}, para
$s\geq r>p-p_0$ existe un monomorfismo canónico
\begin{equation}
\label{0.13.5.5.1-fr}
\tag{13.5.5.1}
i:\mathrm E_s^{pq}(\mathbf X^\bullet)
\longrightarrow\mathrm E_r^{pq}(\mathbf X^\bullet)
\end{equation}
y un isomorfismo canónico
\begin{equation}
\label{0.13.5.5.2-fr}
\tag{13.5.5.2}
j_r:\mathrm E_r^{pq}(\mathbf X^\bullet)
\xrightarrow{\sim}\mathrm E_\infty^{pq}(X_{p+r-1}^\bullet)
\end{equation}
tal que el diagrama
\begin{equation}
\label{0.13.5.5.3-fr}
\tag{13.5.5.3}
\xymatrix{
\mathrm E_s^{pq}(\mathbf X^\bullet)\ar[r]^{j_s}\ar[d]_i
&\mathrm E_\infty^{pq}(X_{p+s-1}^\bullet)\ar[d]\\
\mathrm E_r^{pq}(\mathbf X^\bullet)\ar[r]_{j_r}
&\mathrm E_\infty^{pq}(X_{p+r-1}^\bullet)
}
\end{equation}
es conmutativo, donde la flecha vertical derecha está inducida por el
morfismo
$X_{p+s-1}^\bullet\to X_{p+r-1}^\bullet$.
\end{lemma}

\begin{proof}
La existencia de $i$ proviene de que $\mathrm B_r^{pq}(X_k^\bullet)=\mathrm B_\infty^{pq}(X_k^\bullet)$ para $r>p-p_0$ (\hyperref[0.13.5.3-fr]{13.5.3}, \emph{a)}); se tiene
$\mathrm Z_r^{pq}(X_k^\bullet)=\mathrm Z_\infty^{pq}(X_k^\bullet)$ para $k+1\leq p+r$ (\hyperref[0.13.5.3-fr]{13.5.3}, \emph{b)}), de donde, en particular,
$\mathrm Z_r^{pq}(X_{p+r-1}^\bullet)=
\mathrm Z_\infty^{pq}(X_{p+r-1}^\bullet)$
Por otra parte,
$\mathrm Z_r^{pq}(X_{p+r-1}^\bullet)$ y
$\mathrm B_r^{pq}(X_{p+r-1}^\bullet)$ se identifican canónicamente con
$\mathrm Z_r^{pq}(\mathbf X^\bullet)$ y
$\mathrm B_r^{pq}(\mathbf X^\bullet)$ por
\sref{0.13.5.3},~\emph{c)}.
Esto demuestra la existencia de $j_r$ y la conmutatividad de
\sref{0.13.5.5.3}.
\end{proof}

\begin{corollary}[13.5.6]
\label{0.13.5.6-fr}
Bajo las hipótesis de \textup{(\hyperref[0.13.5.3-fr]{13.5.3})}, si uno de los límites proyectivos $\varprojlim_r\mathrm E_r^{pq}(\mathbf X^\bullet)$, $\varprojlim_k\mathrm E_\infty^{pq}(X_k^\bullet)$ existe, lo mismo ocurre con el otro, y se tiene un isomorfismo canónico
\begin{equation}
\label{0.13.5.6.1-fr}
\tag{13.5.6.1}
j_\infty:\varprojlim_r\mathrm E_r^{pq}(\mathbf X^\bullet)
\xrightarrow{\sim}
\varprojlim_k\mathrm E_\infty^{pq}(X_k^\bullet).
\end{equation}
Además, el sistema proyectivo
$(\mathrm E_r^{pq}(\mathbf X^\bullet))_{r\in\mathbb Z}$ es esencialmente
constante \textup{(\sref{0.13.4.2})} si y solo si el sistema proyectivo
$(\mathrm E_\infty^{pq}(X_k^\bullet))_{k\in\mathbb Z}$ es esencialmente
constante.
\end{corollary}

\begin{env}[13.5.7]
\label{0.13.5.7}
\label{0.13.5.7-fr}
Se denota por
$\mathrm B_\infty^{pq}(\mathbf X^\bullet)$ y
$\mathrm Z_\infty^{pq}(\mathbf X^\bullet)$ los subobjetos de
$\mathrm E_1^{pq}(\mathbf X^\bullet)$ iguales respectivamente a
$\mathrm B_r^{pq}(\mathbf X^\bullet)$ para $r>p-p_0$ y a
$\inf_r\mathrm Z_r^{pq}(\mathbf X^\bullet)$, cuando este último existe.
Así,
$\varprojlim_r\mathrm E_r^{pq}(\mathbf X^\bullet)$
\oldpage[0\textsubscript{III}]{73}
se identifica canónicamente con
$\mathrm E_\infty^{pq}(\mathbf X^\bullet)=
\mathrm Z_\infty^{pq}(\mathbf X^\bullet)/
\mathrm B_\infty^{pq}(\mathbf X^\bullet)$.
Se observará que los objetos
$\mathrm Z_r^{pq}(\mathbf X^\bullet)$,
$\mathrm B_r^{pq}(\mathbf X^\bullet)$,
$\mathrm E_r^{pq}(\mathbf X^\bullet)$
($(1\leq r\leq+\infty)$), y los morfismos $d_r^{pq}$ dependen
\emph{funtorialmente} del sistema proyectivo $\mathbf X^\bullet$, con las
restricciones de \sref{0.13.5.5}, y los morfismos definidos en
\sref{0.13.5.5} y \sref{0.13.5.6} son funtoriales.
\end{env}

\subsection{Sucesión espectral de un funtor relativa a un objeto dotado de una filtración finita}
\label{subsection:0.13.6}
\label{subsection:0.13.6-fr}

\begin{env}[13.6.1]
\label{0.13.6.1}
\label{0.13.6.1-fr}
Sean $\boldsymbol{C}$ y $\boldsymbol{C}'$ dos categorías abelianas, y sea
$\mathrm T:\boldsymbol{C}\to \boldsymbol{C}'$ un funtor covariante
aditivo.
Supongamos que todo objeto de $\boldsymbol{C}$ es isomorfo a un subobjeto de
un objeto inyectivo, de modo que existan los funtores derivados por la derecha
$\mathrm R^p\mathrm T$ ($(p\geq0)$).
\end{env}

\begin{lemma}[13.6.2]
\label{0.13.6.2}
\label{0.13.6.2-fr}
Sea $A$ un objeto de $\boldsymbol{C}$ dotado de una filtración \textup{finita}
$(\mathrm F^i(A))_{i\in\mathbb Z}$.
Existe una resolución inyectiva
$X^\bullet=(X^j)_{j\geq0}$ de $A$, dotada de una filtración finita
$(\mathrm F^i(X^\bullet))_{i\in\mathbb Z}$ tal que la relación $\mathrm F^i(A)=A$ \textup{(resp. $\mathrm F^i(A)=0$)} implique $\mathrm F^i(X^\bullet)=X^\bullet$ \textup{(resp. $\mathrm F^i(X^\bullet)=0$)} y que, para todo $i\in\mathbb Z$,
$\mathrm F^i(X^\bullet)$ es una resolución inyectiva de
$\mathrm F^i(A)$.
\end{lemma}

\begin{proof}
Sea $p$ (resp. $q>p$) el mayor índice tal que
$\mathrm F^i(A)=A$ (resp. el menor índice tal que
$\mathrm F^i(A)=0$).
Se razona por inducción sobre $q-p$, siendo inmediato el lema cuando
$q-p=1$. Habiendo formado una resolución inyectiva ${X''}^\bullet$ de
$A/\mathrm F^{q-1}(A)$ con las propiedades requeridas, se considera la sucesión exacta $0\to\mathrm F^{q-1}(A)\to A\to A/\mathrm F^{q-1}(A)\to0$, se toma una resolución inyectiva ${X'}^\bullet$ de
$\mathrm F^{q-1}(A)$, luego se determina una resolución inyectiva $X^\bullet$ de $A$ de modo que se obtenga una sucesión exacta
$0\to {X'}^\bullet\to X^\bullet\to {X''}^\bullet\to0$
compatible con la sucesión precedente (M, V, 2.2).
Está claro que $X^\bullet$ tiene las propiedades requeridas.
\end{proof}

\begin{corollary}[13.6.3]
\label{0.13.6.3}
\label{0.13.6.3-fr}
Sea $B$ un segundo objeto de $\boldsymbol{C}$, dotado de una filtración
finita $(\mathrm F^i(B))_{i\in\mathbb Z}$, sea $s$ un entero,
y sea $u:A\to B$ un morfismo tal que
$u(\mathrm F^i(A))\subset\mathrm F^{i+s}(B)$
para todo $i\in\mathbb Z$.
Si $Y^\bullet=(Y^j)_{j\geq0}$ es una resolución inyectiva de $B$
dotada de una filtración
$(\mathrm F^i(Y^\bullet))_{i\in\mathbb Z}$ que tenga las propiedades
enunciadas en \sref{0.13.6.2}, entonces existe un morfismo
$v:X^\bullet\to Y^\bullet$, compatible con $u$, tal
que
$v(\mathrm F^i(X^\bullet))\subset\mathrm F^{i+s}(Y^\bullet)$
para todo $i\in\mathbb Z$.
Además, cualesquiera dos de tales morfismos $v$ y $v'$ son homotópicos.
\end{corollary}

Esto se sigue inmediatamente, por inducción sobre $q-p$, de la construcción
precedente y de (M, V, 2.3).

\begin{env}[13.6.4]
\label{0.13.6.4}
\label{0.13.6.4-fr}
Bajo las hipótesis de \sref{0.13.6.2}, consideremos el complejo
$\mathrm T(X^\bullet)$ en $\boldsymbol{C}'$.
Está filtrado por los complejos
$\mathrm T(\mathrm F^i(X^\bullet))$, pues
$\mathrm F^i(X^\bullet)$ es un sumando directo de $X^\bullet$.
Se sigue de \sref{0.13.6.3} que la \emph{sucesión espectral} de
este complejo filtrado depende, salvo isomorfismo, únicamente del
objeto filtrado $A$.
Su término límite es la cohomología $\mathrm R^\bullet\mathrm T(A)$, dotada
de la filtración
\begin{equation}
\label{0.13.6.4.1-fr}
\tag{13.6.4.1}
\begin{split}
\mathrm F^p(\mathrm R^n\mathrm T(A))
&=\operatorname{Im}\bigl(\mathrm R^n\mathrm T(\mathrm F^p(A))
\longrightarrow\mathrm R^n\mathrm T(A)\bigr)\\
&=\operatorname{Ker}\bigl(\mathrm R^n\mathrm T(A)
\longrightarrow\mathrm R^n\mathrm T(A/\mathrm F^p(A))\bigr)
\end{split}
\end{equation}
por \sref{0.11.2.2}, y su término $\mathrm E_1$ es
\begin{equation}
\label{0.13.6.4.2-fr}
\tag{13.6.4.2}
\mathrm E_1^{pq}=\mathrm R^{p+q}\mathrm T(\operatorname{gr}^p(A))
\end{equation}
donde $\operatorname{gr}^p(A)$ designa, como de costumbre, $\mathrm F^p(A)/\mathrm F^{p+1}(A)$. Está claro, por (\hyperref[0.11.2.2-fr]{11.2.2}), que la filtración del término límite es \emph{finita}, y que para $p$ y $q$ dados, las sucesiones de los
\oldpage[0\textsubscript{III}]{74}
$\mathrm B_r^{pq}(A)=\mathrm B_r^{pq}(\mathrm T(X^\bullet))$ y $\mathrm Z_r^{pq}(A)=\mathrm Z_r^{pq}(\mathrm T(X^\bullet))$ son \emph{estacionarias}, por tanto la sucesión espectral precedente es \emph{birregular} (\hyperref[0.11.1.3-fr]{11.1.3}). Denotaremos esta sucesión por $\mathrm E(A)=(\mathrm E_r^{pq}(A))$ y la llamaremos la
\emph{sucesión espectral del funtor $\mathrm T$ relativa al
objeto filtrado $A$}.
\end{env}

\begin{env}[13.6.5]
\label{0.13.6.5}
\label{0.13.6.5-fr}
Supongamos ahora que se cumplen las hipótesis de \sref{0.13.6.3}, y
conservemos su notación.
Como $\mathrm F^i(X^\bullet)$ (resp.
$\mathrm F^i(Y^\bullet)$) es un sumando directo de $X^\bullet$
(resp. $Y^\bullet$), se tiene
$(\mathrm T v)(\mathrm T(\mathrm F^i(X^\bullet)))
\subset\mathrm T(\mathrm F^{i+s}(Y^\bullet))$
para todo $i\in\mathbb Z$.
Las definiciones de \sref{0.11.2.2} muestran entonces que, para
$1\leq r\leq+\infty$, $\mathrm T v$ define morfismos
$\mathrm B_r^{pq}(\mathrm T(X^\bullet))\to
\mathrm B_r^{p+s,q-s}(\mathrm T(Y^\bullet))$
y
$\mathrm Z_r^{pq}(\mathrm T(X^\bullet))\to
\mathrm Z_r^{p+s,q-s}(\mathrm T(Y^\bullet))$,
y, por tanto, un morfismo
\[
w_r:\mathrm E_r^{pq}(A)\to\mathrm E_r^{p+s,q-s}(B)~;
\]
Análogamente, sobre los términos límite se tienen morfismos
$u_n:\mathrm R^n\mathrm T(A)\to\mathrm R^n\mathrm T(B)$
tales que
$u_n(\mathrm F^p(\mathrm R^n\mathrm T(A)))\subset
\mathrm F^{p+s}(\mathrm R^n\mathrm T(B))$.
La definición de los $d_r^{pq}$ (M, XV, 1) muestra también que los
diagramas
\[
\xymatrix{
\mathrm E_r^{pq}(A)\ar[r]^{d_r^{pq}}\ar[d]_{w_r}
&\mathrm E_r^{p+r,q-r+1}(A)\ar[d]^{w_r}\\
\mathrm E_r^{p+s,q-s}(B)\ar[r]_{d_r^{p+s,q-s}}
&\mathrm E_r^{p+r+s,q-r-s+1}(B)
}
\]
son conmutativos.
Se sigue que existe un diagrama conmutativo análogo para los
isomorfismos $\alpha_r^{pq}$, cuya explicitación dejamos al lector.
Finalmente (\emph{loc.\ cit.}), se tienen también diagramas conmutativos
sobre los términos límite:
\[
\xymatrix{
\mathrm E_\infty^{pq}(A)\ar[r]^{\beta^{pq}}\ar[d]_{w_\infty}
&\operatorname{gr}^p(\mathrm R^{p+q}\mathrm T(A))
  \ar[d]^{u_{p+q}}\\
\mathrm E_\infty^{p+s,q-s}(B)\ar[r]_{\beta^{p+s,q-s}}
&\operatorname{gr}^{p+s}(\mathrm R^{p+q}\mathrm T(B))
}
\]
\end{env}

\begin{env}[13.6.6]
\label{0.13.6.6}
\label{0.13.6.6-fr}
Supongamos, en particular, que existe un anillo $S$, dotado de una
\emph{filtración} $(\mathrm F^i(S))_{i\in\mathbb Z}$, y un homomorfismo
de anillos
\begin{equation}
\label{0.13.6.6.1-fr}
\tag{13.6.6.1}
h:S\to\operatorname{Hom}_{\boldsymbol{C}}(A,A)
\end{equation}
\oldpage[0\textsubscript{III}]{75}
tal que, para todo $t\in\mathrm F^j(S)$,
$h_t(\mathrm F^i(A))\subset\mathrm F^{i+j}(A)$ para todo par $i$, $j$. Diremos para abreviar que $A$ está dotado de la estructura de un
\emph{$S$--$\boldsymbol{C}$-módulo filtrado} sobre el anillo filtrado $S$.

Al pasar a los objetos graduados asociados, todo $h_t$, para
$t\in\mathrm F^j(S)$, define un endomorfismo \emph{graduado}
$\overline h_t$ de $\operatorname{gr}^\bullet(A)$, homogéneo de grado $j$; además, este morfismo solo depende de la clase de $t$ en $\operatorname{gr}^j(S)$, y se define así un homomorfismo de \emph{anillos graduados} \[
\overline h:\operatorname{gr}^\bullet(S)\to
\operatorname{Hom}_{\boldsymbol{C}}^g(\operatorname{gr}^\bullet(A),
\operatorname{gr}^\bullet(A))
\] donde el segundo miembro es el anillo de los endomorfismos \emph{graduados} de $\operatorname{gr}^\bullet(A)$.
Decimos que $\operatorname{gr}^\bullet(A)$ está dotado de la estructura
de un \emph{$\operatorname{gr}^\bullet(S)$--$\boldsymbol{C}$-módulo graduado}.

Se sigue de \sref{0.13.6.5} que, para
$1\leq r\leq+\infty$, todo $\overline t\in\operatorname{gr}^j(S)$ define canónicamente en los objetos bigraduados $(\mathrm B_r^{pq}(A))_{(p,q)\in\mathbb Z\times\mathbb Z}$, $(\mathrm Z_r^{pq}(A))_{(p,q)\in\mathbb Z\times\mathbb Z}$ y $\mathrm E_r(A)=(\mathrm E_r^{pq}(A))_{(p,q)\in\mathbb Z\times\mathbb Z}$ endomorfismos bigraduados de grados $(j,-j)$; en $\mathrm E_r(A)$, para $r$ finito, este endomorfismo conmuta con
el endomorfismo bigraduado definido por los $d_r^{pq}$.
Estos endomorfismos satisfacen las condiciones usuales de asociatividad y
distributividad respecto de la adición en
$\operatorname{gr}^\bullet(S)$ y en los objetos bigraduados en cuestión.
Para abreviar, decimos que estos últimos son
\emph{$\operatorname{gr}^\bullet(S)$--$\boldsymbol{C}'$-módulos bigraduados};
es inmediato que los $\alpha_r^{pq}$ definen un \emph{isomorfismo} para
este tipo de estructura.

Para todo entero $n$, sea $\mathrm B_r^{(n)}(A)$
(resp. $\mathrm Z_r^{(n)}(A)$ y
$\mathrm E_r^{(n)}(A)$) el subobjeto graduado de $\mathrm B_r^{\bullet\bullet}(A)$ (resp. $\mathrm Z_r^{\bullet\bullet}(A)$, $\mathrm E_r^{\bullet\bullet}(A)$) formado por los $\mathrm B_r^{pq}(A)$ (resp. $\mathrm Z_r^{pq}(A)$, $\mathrm E_r^{pq}(A)$) tales que $p+q=n$ ($1\leq r\leq+\infty$).
Es inmediato que estos son
\emph{$\operatorname{gr}^\bullet(S)$--$\boldsymbol{C}'$-módulos graduados}.

Finalmente, todo $\overline t\in\operatorname{gr}^j(S)$ define, para todo
$n$ un endomorfismo graduado de grado $j$ en el objeto graduado $\operatorname{gr}^\bullet(\mathrm R^n\mathrm T(A))$, que adquiere así la
estructura de un
$\operatorname{gr}^\bullet(S)$--\emph{$\boldsymbol{C}'$-módulo graduado}; los $\beta^{pq}$ (para $p+q=n$) definen un \emph{isomorfismo} de $\mathrm E_\infty^{(n)}(A)$ sobre $\operatorname{gr}^\bullet(\mathrm R^n\mathrm T(A))$ para esta especie de estructura.

Se observará que cuando $\boldsymbol{C}'$ es la categoría de \emph{grupos abelianos}, las
estructuras de $S$--$\boldsymbol{C}'$-módulo (resp. de
$\operatorname{gr}^\bullet(S)$--$\boldsymbol{C}'$-módulo graduado o bigraduado) son
simplemente las estructuras usuales de $S$-\emph{módulo} (resp. de
$\operatorname{gr}^\bullet(S)$-módulo graduado o bigraduado).
\end{env}

\subsection{Funtores derivados de un límite proyectivo de argumentos}
\label{subsection:0.13.7}
\label{subsection:0.13.7-fr}
\label{0.13.7}

\begin{env}[13.7.1]
\label{0.13.7.1}
\label{0.13.7.1-fr}
Sean $\boldsymbol{C}$ y $\boldsymbol{C}'$ dos categorías abelianas, suponiéndose $\boldsymbol{C}$ tal que todo objeto de $\boldsymbol{C}$ es isomorfo a un subobjeto de un
objeto inyectivo.
Sea
$\mathrm T:\boldsymbol{C}\to \boldsymbol{C}'$
un funtor covariante aditivo.
Consideremos un \emph{sistema proyectivo estricto}
$\mathbf A=(A_k)_{k\in\mathbb Z}$ en $\boldsymbol{C}$, acotado inferiormente; más
precisamente, supongamos que $A_k=0$ para $k<k_0$. Asociamos canónicamente a este sistema una filtración $(\mathrm F^p(A_k))_{p\in\mathbb Z}$ sobre cada $A_k$ mediante las fórmulas (\hyperref[0.13.4.1.1-fr]{13.4.1.1}), y como se trata de un sistema proyectivo \emph{estricto}, los morfismos canónicos
\begin{equation}
\label{0.13.7.1.1-fr}
\tag{13.7.1.1}
\mathrm F^i(A_h)/\mathrm F^j(A_h)\to
\mathrm F^i(A_k)/\mathrm F^j(A_k)\qquad(h\geq k)
\end{equation}
son \emph{isomorfismos} para $i\leq j\leq k+1$ son \emph{isomorfismos}. Recordemos además que se tiene $\mathrm F^{k_0}(A_k)=A_k$ y
$\mathrm F^{k+1}(A_k)=0$ para todo $k$.
\end{env}

\begin{env}[13.7.2]
\label{0.13.7.2}
\label{0.13.7.2-fr}
Para cada $k$ una \emph{resolución inyectiva} $X_k^\bullet=(X_k^j)_{j\geq0}$ de $A_k$ que tenga las
propiedades de \sref{0.13.6.2}.
Los morfismos canónicos $A_{k+1}\to A_k$ nos permiten, por
\sref{0.13.6.3}, definir para cada $k$ un morfismo de complejos
$X_{k+1}^\bullet\to X_k^\bullet$ compatibles
\oldpage[0\textsubscript{III}]{76}
con las filtraciones, y que hacen de
$\mathbf X^\bullet=(X_k^\bullet)_{k\in\mathbb Z}$ un \emph{sistema proyectivo} de complejos. Se puede suponer además que este sistema proyectivo es \emph{estricto}. Para ello, se observa que, en virtud del isomorfismo (\hyperref[0.13.7.1.1-fr]{13.7.1.1}), $A_k$ es isomorfo a $A_{k+1}/\mathrm F^{k+1}(A_{k+1})$; se puede, pues, tomar, en la construcción de $X_{k+1}^\bullet$, la resolución inyectiva de $A_{k+1}/\mathrm F^{k+1}(A_{k+1})$ \emph{igual} a
$X_k^\bullet$.
Se sigue entonces de (M, V, 2.3) que el morfismo de complejos
$X_{k+1}^\bullet\to X_k^\bullet$ puede
construirse de modo que, al pasar a los cocientes, induzca un
\emph{isomorfismo}
$X_{k+1}^\bullet/\mathrm F^{k+1}(X_{k+1}^\bullet)
\xrightarrow{\sim}X_k^\bullet$
que respeta las filtraciones; esta es precisamente la condición de
\sref{0.13.5.1}.
\end{env}

\begin{env}[13.7.3]
\label{0.13.7.3}
\label{0.13.7.3-fr}
Por construcción, el sistema proyectivo
$\mathrm T(\mathbf X^\bullet)$ de complejos
$\mathrm T(X_k^\bullet)$ satisface las hipótesis de
\sref{0.13.5.3}.
Los resultados de \sref{0.13.5.4}, \sref{0.13.5.5} y
\sref{0.13.5.6} se aplican, por tanto, a las sucesiones espectrales
$\mathrm E(\mathrm T(X_k^\bullet))=\mathrm E(A_k)$; escribiremos $\mathrm E_r^{pq}(\mathbf A)$ en lugar de
$\mathrm E_r^{pq}(\mathrm T(\mathbf X^\bullet))$ para $1\leq r\leq+\infty$ (cf. (\hyperref[0.13.5.7-fr]{13.5.7}) para $r=+\infty$) y lo mismo para las notaciones análogas. Se observará en particular que se tiene \begin{equation}
\label{0.13.7.3.1-fr}
\tag{13.7.3.1}
\mathrm E_1^{pq}(\mathbf A)=
\mathrm R^{p+q}\mathrm T(\operatorname{gr}^p(\mathbf A))
\end{equation}
por \sref{0.13.6.4.2} y el hecho de que el sistema proyectivo
$(\operatorname{gr}^p(A_k))$ es esencialmente constante.
Estos resultados, junto con \sref{0.13.4.3}, dan la proposición
siguiente, demostrada por primera vez por Shih Weishu mediante un método
diferente (inédito).
\end{env}

\begin{proposition}[13.7.4]
\label{0.13.7.4}
\label{0.13.7.4-fr}
(Shih).
Sea $n$ un entero.
Las dos condiciones siguientes son equivalentes:
\begin{enumerate}
  \item[{\rm a)}] Para todo par $(p,q)$ tal que $p+q=n$, el
  sistema proyectivo
  $(\mathrm E_r^{pq}(\mathbf A))_{r\geq2}$ es esencialmente constante.
  \item[{\rm b)}] El sistema proyectivo
  $\mathrm R^n\mathrm T(\mathbf A)=
(\mathrm R^n\mathrm T(A_k))_{k\in\mathbb Z}$ satisface \textup{(ML)}.
\end{enumerate}
Además, cuando se cumplen estas condiciones, se tiene un isomorfismo canónico
\begin{equation}
\label{0.13.7.4.1-fr}
\tag{13.7.4.1}
\operatorname{gr}^p(\mathrm R^n\mathrm T(\mathbf A))
\xrightarrow{\sim}\mathrm E_\infty^{p,n-p}(\mathbf A)
\qquad\textit{pour tout $p\in\mathbb Z$.}
\end{equation}
\end{proposition}
\begin{proof}
En efecto, en virtud de (\hyperref[0.13.5.6-fr]{13.5.6}), la condición \emph{a)} equivale a decir que el sistema proyectivo $(\mathrm E_\infty^{pq}(A_k))_{k\in\mathbb Z}$ es esencialmente constante para $p+q=n$, y por otra parte
$\operatorname{gr}^p(\mathrm R^n\mathrm T(A_k))$ es canónicamente isomorfo a
$\mathrm E_\infty^{p,n-p}(A_k)$.
Se sigue de \sref{0.13.4.3} que {\rm a)} y {\rm b)} son
equivalentes, y \sref{0.13.7.4.1} no es sino
\sref{0.13.5.6.1} en el caso considerado.
\end{proof}

\begin{corollary}[13.7.5]
\label{0.13.7.5}
\label{0.13.7.5-fr}
Sea $\boldsymbol{\mathcal F}=(\mathcal F_k)_{k\in\mathbb N}$ un sistema proyectivo de
haces de grupos abelianos que satisface las condiciones {\rm(i)}, {\rm(ii)}
y {\rm(iii)} de \sref{0.13.3.1}, y sea
$\mathcal F=\varprojlim\mathcal F_k$.
Supongamos que, para el funtor
$\mathcal G\leadsto\Gamma(X,\mathcal G)$, el sistema proyectivo $(\mathrm E_r^{pq}(\boldsymbol{\mathcal F}))_{r\in\mathbb Z}$ sea esencialmente constante para todo par $(p,q)$ tal que $p+q=n$ o $p+q=n+1$. Consideremos sobre $H^{n+1}(X,\mathcal F)$ la filtración definida por $\mathrm F^p(H^{n+1}(X,\mathcal F))=
\operatorname{Ker}(H^{n+1}(X,\mathcal F)\to
H^{n+1}(X,\mathcal F_{p-1}))$.
Entonces existe un isomorfismo canónico
\begin{equation}
\label{0.13.7.5.1-fr}
\tag{13.7.5.1}
\operatorname{gr}^p(H^{n+1}(X,\mathcal F))
\xrightarrow{\sim}\mathrm E_\infty^{p,n-p+1}(\boldsymbol{\mathcal F})
\qquad\textit{pour tout $p\in\mathbb Z$.}
\end{equation}
\end{corollary}

\begin{proof}
Resulta de (\hyperref[0.13.7.4-fr]{13.7.4}) aplicado al caso en que $\boldsymbol{C}$ la categoría de haces de
grupos abelianos sobre $X$, como $\boldsymbol{C}'$ la categoría de grupos abelianos,
y $\mathrm T=\Gamma$, que se tiene un isomorfismo
\oldpage[0\textsubscript{III}]{77}
canónico
$\operatorname{gr}^p(\mathrm R^{n+1}\Gamma(\boldsymbol{\mathcal F}))
\xrightarrow{\sim}\mathrm E_\infty^{p,n-p+1}(\boldsymbol{\mathcal F})$ para todo $p\in\mathbb Z$. Por otra parte, como, en virtud de (\hyperref[0.13.7.4-fr]{13.7.4}), el sistema proyectivo $(H^n(X,\mathcal F_k))_{k\in\mathbb Z}$ satisface (ML), se deduce de (\hyperref[0.13.3.1-fr]{13.3.1}) un isomorfismo canónico
\hypertarget{canon-legacy.0-III.13.7.5.proof-comparison}{}
\[
\label{0.13.7.5.1b-fr}
\tag{13.7.5.2}
H^{n+1}(X,\mathcal F)\xrightarrow{\sim}
\varprojlim_k H^{n+1}(X,\mathcal F_k).
\]

Como el sistema proyectivo
$\mathrm R^{n+1}\Gamma(\boldsymbol{\mathcal F})$ satisface {\rm(ML)} por
\sref{0.13.7.4}, existe un isomorfismo canónico
$\operatorname{gr}^p(\mathrm R^{n+1}\Gamma(\boldsymbol{\mathcal F}))
\xrightarrow{\sim}
\varprojlim_k\operatorname{gr}^p(H^{n+1}(X,\mathcal F_k))$
por \sref{0.13.4.3}, y un isomorfismo canónico
$\varprojlim_k\operatorname{gr}^p(H^{n+1}(X,\mathcal F_k))
\xrightarrow{\sim}
\operatorname{gr}^p(\varprojlim_k H^{n+1}(X,\mathcal F_k))$
(\hyperref[0.13.4.5-fr]{13.4.5}). Todo se reduce, pues, a ver que el isomorfismo (\hyperref[0.13.7.5.1b-fr]{13.7.5.2}) es compatible con las filtraciones de los dos miembros; pero esto resulta inmediatamente de las definiciones y de la conmutatividad del diagrama
\[
\xymatrix{
H^{n+1}(X,\mathcal F)\ar[rr]^{\sim}\ar[dr]
&&\varprojlim_k H^{n+1}(X,\mathcal F_k)\ar[dl]\\
&H^{n+1}(X,\mathcal F_{p-1})&
}
\]
para todo $p$.
\end{proof}

\begin{env}[13.7.6]
\label{0.13.7.6}
\label{0.13.7.6-fr}
Sea $S$ un anillo dotado de una filtración
$(\mathrm F^i(S))_{i\in\mathbb Z}$ tal que
$\mathrm F^0(S)=S$ (y por tanto
$\operatorname{gr}^i(S)=0$ para $i<0$).
Supongamos que cada $A_k$, con la filtración definida en
\sref{0.13.7.1}, es un
$S$--$\boldsymbol{C}$-\emph{módulo filtrado} en el sentido de \sref{0.13.6.6}, y que los
morfismos $A_h\to A_k$ para $k\leq h$ son morfismos de
$S$--$\boldsymbol{C}$-módulos filtrados.
Para abreviar, diremos entonces que $\mathbf A$ es un
\emph{sistema proyectivo de $S$--$\boldsymbol{C}$-módulos filtrados}. Entonces es inmediato que los morfismos $\mathrm B_r^{pq}(A_h)\to\mathrm B_r^{pq}(A_k)$ y $\mathrm Z_r^{pq}(A_h)\to\mathrm Z_r^{pq}(A_k)$ para $k\leq h$ y
$1\leq r\leq+\infty$, son morfismos para las estructuras de $\operatorname{gr}^\bullet(S)$-$\boldsymbol{C}'$-\emph{módulo bigraduado} (\hyperref[0.13.6.5-fr]{13.6.5}), y que las familias $(\mathrm Z_r^{pq}(\mathbf A))$, $(\mathrm B_r^{pq}(\mathbf A))$ y $(\mathrm E_r^{pq}(\mathbf A))$ son $\operatorname{gr}^\bullet(S)$--$\boldsymbol{C}'$-\emph{módulos bigraduados}
(\sref{0.13.6.5}), y que, para $r$ finito, siendo los dos primeros submódulos de $(\mathrm E_1^{pq}(\mathbf A))$. Se denotará también por $\mathrm Z_r^{(n)}(\mathbf A)$,
$\mathrm B_r^{(n)}(\mathbf A)$ y
$\mathrm E_r^{(n)}(\mathbf A)$ para los respectivos subobjetos graduados
obtenidos conservando únicamente los términos con $p+q=n$; estos son
$\operatorname{gr}^\bullet(S)$--$\boldsymbol{C}'$-\emph{módulos graduados}.

Cuando el sistema proyectivo
$(\mathrm E_r^{pq}(\mathbf A))_{r\in\mathbb Z}$ es esencialmente constante,
$(\mathrm E_\infty^{pq}(\mathbf A))$ es asimismo un
$\operatorname{gr}^\bullet(S)$--$\boldsymbol{C}'$-módulo bigraduado, y cada
$\mathrm E_\infty^{(n)}(\mathbf A)$ es un
$\operatorname{gr}^\bullet(S)$--$\boldsymbol{C}'$-módulo graduado. Además, los $\beta^{p,n-p}:\mathrm E_\infty^{p,n-p}(A_k)\xrightarrow{\sim}
\operatorname{gr}^p(\mathrm R^n\mathrm T(A_k))$ constituyen para cada $k$ un isomorfismo para la estructura de $\operatorname{gr}^\bullet(S)$--$\boldsymbol{C}'$-módulos graduados de
$\mathrm E_\infty^{(n)}(A_k)$ sobre
$\operatorname{gr}^\bullet(\mathrm R^n\mathrm T(A_k))$.
Bajo las hipótesis precedentes, el isomorfismo
$\beta^{p,n-p}:\mathrm E_\infty^{(n)}(\mathbf A)\xrightarrow{\sim}
\varprojlim_k\operatorname{gr}^\bullet(\mathrm R^n\mathrm T(A_k))$
respeta, por tanto, también estas estructuras.
Lo mismo vale manifiestamente para el isomorfismo canónico
$\operatorname{gr}^\bullet(\mathrm R^n\mathrm T(\mathbf A))
\xrightarrow{\sim}
\varprojlim_k\operatorname{gr}^\bullet(\mathrm R^n\mathrm T(A_k))$,
por consiguiente, los isomorfismos \sref{0.13.7.4.1} forman un
isomorfismo de
$\operatorname{gr}^\bullet(S)$--$\boldsymbol{C}'$-\emph{módulos graduados}.
\end{env}

\begin{proposition}[13.7.7]
\label{0.13.7.7}
\label{0.13.7.7-fr}
Sea $S$ un anillo noetheriano $\mathfrak J$-ádico.
Supongamos que $\boldsymbol{C}$ es una categoría abeliana en la que todo objeto
es isomorfo a un subobjeto de un objeto inyectivo, y sea
$\mathrm T$ un funtor covariante aditivo de $\boldsymbol{C}$ en la
categoría de grupos abelianos.
Sea $\mathbf A=(A_k)_{k\in\mathbb Z}$ un
\oldpage[0\textsubscript{III}]{78}
sistema proyectivo estricto de $S$--$\boldsymbol{C}$-módulos filtrados, para la
filtración $\mathfrak J$-ádica sobre $S$, acotado inferiormente.
Para un entero fijo $n$ dado, se satisfaga la condición siguiente:
\begin{itemize}
\item[$(\mathrm F_n)$]
  \emph{El $\operatorname{gr}^\bullet(S)$-módulo graduado}
  $\mathrm E_1^{(m)}(\mathbf A)=
(\mathrm R^m\mathrm T(\operatorname{gr}^p(\mathbf A)))_{p\in\mathbb Z}$
  \emph{de \sref{0.13.7.3.1} es de tipo finito para
  $m=n$ y $m=n+1$.}
\end{itemize}
Entonces:
\begin{enumerate}
  \item[{\rm(i)}] Los sistemas proyectivos
  $(\mathrm R^n\mathrm T(A_k))_{k\in\mathbb Z}$ y
  $(\mathrm R^{n+1}\mathrm T(A_k))_{k\in\mathbb Z}$ satisfacen {\rm(ML)}.
  \item[{\rm(ii)}] Si
  $\mathrm R'^{\,n}\mathrm T(\mathbf A)=
\varprojlim_k\mathrm R^n\mathrm T(A_k)$,
  entonces $\mathrm R'^{\,n}\mathrm T(\mathbf A)$ es un $S$-módulo de tipo finito.
\item[\rm(iii)] La filtración definida por $\mathrm F^p(\mathrm R'^{\,n}\mathrm T(\mathbf A))=
\operatorname{Ker}(\mathrm R'^{\,n}\mathrm T(\mathbf A)
\to\mathrm R^n\mathrm T(A_{p-1}))$ $(p\in\mathbb Z)$ sobre $\mathrm R'^{\,n}\mathrm T(\mathbf A)$ es $\mathfrak J$-buena; es decir,
  $\mathfrak J\mathrm F^p(\mathrm R'^{\,n}\mathrm T(\mathbf A))
\subset\mathrm F^{p+1}(\mathrm R'^{\,n}\mathrm T(\mathbf A))$
  para todo $p$, con igualdad para todo $p$ suficientemente grande.
  En particular, la topología sobre
  $\mathrm R'^{\,n}\mathrm T(\mathbf A)$ definida por esta filtración es la
  topología $\mathfrak J$-ádica.
  \item[{\rm(iv)}] El sistema proyectivo
  $(\mathrm E_r^{pq}(\mathbf A))_{r\in\mathbb Z}$ es esencialmente constante siempre que
  $p+q=n$ o $p+q=n+1$.
  Así, $\mathrm E_\infty^{pq}(\mathbf A)$ está, pues, definido \textup{(\hyperref[0.13.5.7-fr]{13.5.7})} y se tiene un isomorfismo canónico de $\operatorname{gr}^\bullet(S)$-módulos graduados
\begin{equation}
\label{0.13.7.7.1-fr}
\tag{13.7.7.1}
\operatorname{gr}^p(\mathrm R'^{\,n}\mathrm T(\mathbf A))
\xrightarrow{\sim}\mathrm E_\infty^{p,n-p}(\mathbf A)
\quad(p\in\mathbb Z).
\end{equation}
\end{enumerate}
\end{proposition}

\begin{proof}
Se observará que el isomorfismo (\hyperref[0.13.7.7.1-fr]{13.7.7.1}) permitirá denotar por $\mathrm R^n\mathrm T(\mathbf A)$ para el límite proyectivo
$\mathrm R'^{\,n}\mathrm T(\mathbf A)$ del sistema proyectivo
$\mathrm R^n\mathrm T(\mathbf A)$.

El anillo graduado $\operatorname{gr}^\bullet(S)$ es noetheriano
(Bourbaki, \emph{Alg.\ comm.}, cap.~III, \textsection2,
n\textsuperscript{o}~9, cor.~5 del teorema~2).
Por tanto, la sucesión creciente de
$\operatorname{gr}^\bullet(S)$-submódulos graduados
$\mathrm B_r^{(m)}(\mathbf A)$ de
$\mathrm E_1^{(m)}(\mathbf A)$ es \emph{estacionaria} para $m=n$ y
$m=n+1$.
La condición \emph{(b)} de \sref{0.11.1.10} queda, por tanto, satisfecha.
Se sigue que la condición \emph{(a)} de \sref{0.13.7.4} se cumple para
$n$ y para $n+1$, lo que demuestra {\rm(i)}.

Además, los isomorfismos \sref{0.13.7.4.1}, junto con las
observaciones de \sref{0.13.7.6}, muestran que
$\operatorname{gr}^\bullet(\mathrm R^n\mathrm T(\mathbf A))$ es un
$\operatorname{gr}^\bullet(S)$-módulo graduado isomorfo a
$\mathrm E_\infty^{(n)}(\mathbf A)=
\mathrm Z_\infty^{(n)}(\mathbf A)/\mathrm B_\infty^{(n)}(\mathbf A)$
Como $\mathrm Z_\infty^{(n)}(\mathbf A)$ es un submódulo de
$\mathrm E_1^{(n)}(\mathbf A)$, es de \emph{tipo finito}, y también lo es
$\mathrm E_\infty^{(n)}(\mathbf A)$.
Finalmente, para la filtración
$(\mathrm F^p(\mathrm R'^{\,n}\mathrm T(\mathbf A)))$, resulta de (\hyperref[0.13.4.5-fr]{13.4.5}) que $\operatorname{gr}^\bullet(\mathrm R^n\mathrm T(\mathbf A))$ y $\operatorname{gr}^\bullet(\mathrm R'^{\,n}\mathrm T(\mathbf A))$ son $\operatorname{gr}^\bullet(S)$-módulos graduados isomorfos.
Esto demuestra {\rm(iv)}.
Las aserciones {\rm(ii)} y {\rm(iii)} se seguirán de estos resultados
y del lema siguiente.
\end{proof}

\begin{lemma}[13.7.7.2]
\label{0.13.7.7.2}
\label{0.13.7.7.2-fr}
Sea $S$ un anillo noetheriano $\mathfrak J$-ádico, y sea $M$
un $S$-módulo dotado de una filtración codiscreta
$(\mathrm F^p(M))_{p\in\mathbb Z}$ tal que
$\mathfrak J\mathrm F^p(M)\subset\mathrm F^{p+1}(M)$
Así, $M$ es un módulo filtrado sobre $S$ dotado de su
filtración $\mathfrak J$-ádica.
Supongamos además que $M$ es separado para la topología definida por
$(\mathrm F^p(M))$.
Entonces las condiciones siguientes son equivalentes:
\begin{enumerate}
  \item[{\rm a)}] $M$ es un $S$-módulo de tipo finito y
  $(\mathrm F^p(M))$ es una filtración $\mathfrak J$-buena.
  \item[{\rm b)}] $\operatorname{gr}^\bullet(M)$ es un
  $\operatorname{gr}^\bullet(S)$-módulo de tipo finito.
  \item[{\rm c)}] Los $\operatorname{gr}^p(M)$ son $S$-módulos de
  tipo finito y, para todo $p$ suficientemente grande, los homomorfismos
  canónicos
  \begin{equation}
\label{0.13.7.7.3-fr}
\tag{13.7.7.3}
\mathfrak J\otimes_S\operatorname{gr}^p(M)\longrightarrow
\operatorname{gr}^{p+1}(M)
\end{equation}
  \oldpage[0\textsubscript{III}]{79}
  son sobreyectivos.
  Están inducidos por
  $\mathfrak J\otimes_S\mathrm F^p(M)\to\mathrm F^{p+1}(M)$, teniendo en cuenta el hecho
  de que la imagen del homomorfismo compuesto
  $\mathfrak J\otimes_S\mathrm F^{p+1}(M)\to
\mathfrak J\otimes_S\mathrm F^p(M)\to\mathrm F^{p+1}(M)$
  es
  $\mathfrak J\mathrm F^{p+1}(M)\subset\mathrm F^{p+2}(M)$.
\end{enumerate}
\end{lemma}

Para la demostración, véase Bourbaki, \emph{Alg.\ comm.}, cap.~III,
\textsection3, n\textsuperscript{o}~1, prop.~3.

\begin{env}[13.7.7.4]
\label{0.13.7.7.4}
\label{0.13.7.7.4-fr}
Para aplicar el lema \sref{0.13.7.7.2}, queda observar que la
topología sobre $\mathrm R'^{\,n}\mathrm T(\mathbf A)$ por la filtración considerada hace de $\mathrm R'^{\,n}\mathrm T(\mathbf A)$ un $S$-módulo separado y completo:
esta es la topología de límite proyectivo de los grupos discretos
$\mathrm R^n\mathrm T(A_k)$.
Por otra parte, si $A_k=0$ para $k<k_0$, entonces también
$\mathrm R^n\mathrm T(A_k)=0$ para $k<k_0$; por tanto
$\mathrm F^{k_0}(\mathrm R'^{\,n}\mathrm T(\mathbf A))=
\mathrm R'^{\,n}\mathrm T(\mathbf A)$,
Así, todas las hipótesis del lema quedan efectivamente satisfechas.
\end{env}

\begin{corollary}[13.7.8]
\label{0.13.7.8}
\label{0.13.7.8-fr}
Si se cumple la hipótesis $(\mathrm F_n)$, entonces para todo $f\in S$
existe un isomorfismo canónico
\begin{equation}
\label{0.13.7.8.1-fr}
\tag{13.7.8.1}
\varprojlim_k\bigl((\mathrm R^n\mathrm T(A_k))_f\bigr)
\xrightarrow{\sim}\mathrm R^n\mathrm T(\mathbf A)\otimes_S S_{\{f\}}.
\end{equation}
\end{corollary}

\begin{proof}
En efecto, $\mathrm R^n\mathrm T(\mathbf A)$ es un $S$-módulo de tipo finito,
y $S_{\{f\}}$ es una $S$-álgebra ádica noetheriana
(\sref[0\textsubscript{I}]{0.7.6.11}), a saber, la completación separada
de $S_f$ para la topología $\mathfrak J$-preádica
(\sref[0\textsubscript{I}]{0.7.6.2}).
Se sigue de \sref[0\textsubscript{I}]{0.7.7.8} y
\sref[0\textsubscript{I}]{0.7.7.1} que
$\mathrm R^n\mathrm T(\mathbf A)\otimes_S S_{\{f\}}$
es isomorfo a la completación separada de
$\mathrm R^n\mathrm T(\mathbf A)\otimes_S S_f$ para la
topología $\mathfrak J$-preádica; un sistema fundamental de vecindades de $0$ para esta topología es
$(\mathfrak J^p\mathrm R^n\mathrm T(\mathbf A))\otimes_S S_f$,
los grupos
$\mathrm F^p(\mathrm R^n\mathrm T(\mathbf A))\otimes_S S_f$
forman, por tanto, también un sistema tal.
Este último grupo es el núcleo de la aplicación canónica
$(\mathrm R^n\mathrm T(\mathbf A))_f\to
(\mathrm R^n\mathrm T(A_{p-1}))_f$
Por consiguiente, el grupo separado asociado a
$\mathrm R^n\mathrm T(\mathbf A)\otimes_S S_f$ se identifica con un subgrupo
$G$ de
$\varprojlim_k((\mathrm R^n\mathrm T(A_k))_f)$.
Pero el sistema proyectivo
$((\mathrm R^n\mathrm T(A_k))_f)$ satisface manifiestamente {\rm(ML)}, y la
imagen de $(\mathrm R^n\mathrm T(\mathbf A))_f$ en cada
$(\mathrm R^n\mathrm T(A_k))_f$ es igual a la imagen común de los
$(\mathrm R^n\mathrm T(A_h))_f$ para todo $h\geq k$ suficientemente grande.
Se sigue inmediatamente que $G$ es \emph{denso en todas partes} en
$\varprojlim_k((\mathrm R^n\mathrm T(A_k))_f)$.
Como este último grupo es separado y completo, queda demostrado el
corolario.
\end{proof}
